\chapter{Đào sâu: Gateway dưới tải --- bên trong stack reactive, rate limit và việc của biên}
\label{ch:dd-gateway}

Chương~4 dạy gateway như một bảng route: predicate, URI \code{lb://},
\code{StripPrefix}, một filter JWT. Đó là gateway vào một ngày yên ả.
Chương này là gateway vào một ngày tệ --- course-service chậm, một
client kẹt trong vòng lặp retry, một pod khởi động lại đúng lúc không
nên --- và hỏi stack reactive thực sự mua được gì, biên của dự án đang
thiếu gì, và ranh giới giữa việc của biên và việc của mọi người khác
nằm ở đâu.

Ba câu hỏi sắp xếp chương này:

\begin{enumerate}
  \item \textbf{Vì sao gateway không chặn (non-blocking)}, và chính
  xác chuyện gì hỏng nếu một filter chặn?
  \item \textbf{Cái gì gãy trước dưới tải} --- và vì sao câu trả lời là
  ``thứ không có timeout''?
  \item \textbf{Cái gì thuộc về biên} --- rate limit, circuit breaker,
  CORS --- và cái gì không bao giờ nên?
\end{enumerate}

\section{Vì sao cần gateway, và nó giải quyết gì}

Không có gateway, frontend phải giữ sáu địa chỉ, mỗi service phải tự
xác minh JWT, và CORS, rate limit, header truy vết sẽ được làm sáu lần
hoặc không lần nào. Gateway gom tất cả về một địa chỉ và một chỗ cho
mọi mối quan tâm \emph{xuyên suốt} --- những thứ thuộc về \emph{mọi}
request chứ không thuộc về khóa học hay người dùng. Dạng phỏng vấn của
câu hỏi là ``API gateway làm được gì mà load balancer không làm được?''
và câu trả lời senior là: load balancer chọn một backend; gateway
\emph{hiểu request} --- nó có thể xác thực, viết lại, giới hạn tốc độ
và chuyển hướng theo từng route, vì nó chạy code ứng dụng trên đường
đi. Sức mạnh đó cũng là rủi ro: bất cứ thứ gì bạn đặt lên đường đó, mọi
request đều phải trả giá.

\begin{decision}[gateway ứng dụng phía sau router hạ tầng]
Cụm đã có Traefik (Chương~18). Vì sao vẫn giữ Spring Cloud Gateway
phía sau? Vì hai thứ làm hai việc khác nhau: Traefik sở hữu host, TLS
và việc chia nhánh theo path; gateway sở hữu những thứ chỉ ứng dụng
mới biết --- tên Eureka nào phục vụ \code{/api/courses}, JWT hợp lệ
trông ra sao, header nào cần chèn. Gộp chúng nghĩa là dạy Traefik
secret JWT và hợp đồng header, việc khả thi (middleware ForwardAuth)
nhưng đưa bảo mật ứng dụng vào cấu hình hạ tầng. Quyết định sẽ lật nếu
nền tảng tự cung cấp danh tính (API Gateway cloud với Cognito hay IAP):
lúc đó gateway Spring là một chặng thứ hai trả tiền mà không được gì.
\end{decision}

\section{Stack reactive, đọc từ filter}
\label{sec:reactive}

Gateway khai báo \code{web-application-type: reactive} trong
\code{application.yml} của nó. Một dòng đó đổi mô hình mỗi request một
thread của Tomcat lấy vòng lặp sự kiện (event loop) của Reactor Netty,
và mọi dòng của filter JWT được viết cho mô hình đó:

\begin{lstlisting}[language=Java]
@Override
public Mono<Void> filter(ServerWebExchange exchange,
                         GatewayFilterChain chain) {           (*@\co{1}@*)
    String path = exchange.getRequest().getURI().getPath();
    if (isPublicPath(path)) {
        return chain.filter(exchange);                          (*@\co{2}@*)
    }
    String authHeader = exchange.getRequest().getHeaders()
            .getFirst(HttpHeaders.AUTHORIZATION);
    if (authHeader == null || !authHeader.startsWith("Bearer ")) {
        return unauthorizedResponse(exchange, "Missing ...");  (*@\co{3}@*)
    }
    String token = authHeader.substring(7);
    try {
        Claims claims = Jwts.parser().verifyWith(signingKey)
                .build().parseSignedClaims(token).getPayload(); (*@\co{4}@*)
        ServerHttpRequest mutated = exchange.getRequest().mutate()
                .header(JwtConstants.HEADER_USER_ID, claims.getSubject())
                .header(JwtConstants.HEADER_USER_ROLE,
                        claims.get(JwtConstants.CLAIM_ROLE, String.class))
                .build();
        return chain.filter(exchange.mutate().request(mutated).build());
    } catch (Exception e) {
        return unauthorizedResponse(exchange, "Invalid or expired token");
    }
}
\end{lstlisting}

\begin{description}
  \item[\co{1}] Kiểu trả về là toàn bộ câu chuyện. Một filter servlet
  trả về \code{void} \emph{sau khi} request đã xong; cái này trả về một
  \code{Mono<Void>} --- một \emph{mô tả} công việc sẽ hoàn thành sau.
  Thread gọi \code{filter} được tự do ngay khi phương thức này trả về.
  \item[\co{2}] \code{chain.filter(exchange)} không chuyển tiếp request;
  nó trả về \code{Mono} của chặng kế tiếp. Chưa có gì xảy ra. Chuỗi
  được lắp ráp trước và subscribe một lần.
  \item[\co{3}] Ngay cả nhánh lỗi cũng reactive: \code{writeWith} trả
  về một \code{Mono} hoàn thành khi Netty đã flush xong body.
  \item[\co{4}] Mẩu công việc thật duy nhất, và nó thuộc loại tốn CPU
  chỉ vài micro giây: một HMAC trên vài trăm byte. Đây là loại việc
  \emph{được phép} chạy trên thread event loop. Một lời gọi cơ sở dữ
  liệu thì không.
\end{description}

\subsection{Thread: bốn, không phải bốn trăm}

Tomcat trả lời tính đồng thời bằng thread: mặc định 200, mỗi request
đang bay một thread, mỗi thread đứng chờ trong lúc backend trả lời.
Request của một proxy gần như toàn là chờ, nên một proxy Tomcat với
1.000 request đang bay cần 1.000 thread và khoảng một gigabyte cho
stack. Reactor Netty chạy một nhóm nhỏ, cố định các thread
\emph{event loop}, mỗi thread ghép kênh hàng nghìn socket bằng I/O
không chặn: một request đến, chuỗi filter được lắp ráp, lệnh ghi lên
upstream được phát đi, và thread đi tiếp. Khi upstream trả lời, vòng
lặp nào sở hữu socket đó sẽ tiếp tục chuỗi.

\begin{hardware}
Reactor Netty định cỡ vòng lặp bằng \code{max(availableProcessors, 4)}.
Trên hai lõi của ts-server, đó là \textbf{bốn} thread event loop cho
toàn bộ gateway. Bốn thread phục vụ mọi request là lý do Deployment của
gateway xin ít bộ nhớ hơn một service MVC (384\,Mi so với 512\,Mi,
Chương~17) --- và là lý do một thread bị chặn không phải vấn đề
0,5\,\% mà là vấn đề 25\,\%.
\end{hardware}

\subsection{Chặn làm gì một event loop}

Giả sử ai đó ``cải tiến'' filter bằng cách kiểm tra danh sách cấm trong
Postgres:

\begin{lstlisting}[language=Java]
// DO NOT: a JDBC call inside a WebFlux filter
boolean revoked = jdbcTemplate.queryForObject(
        "select exists(select 1 from revoked where jti = ?)",
        Boolean.class, claims.getId());
\end{lstlisting}

JDBC chặn theo đặc tả: thread gọi đứng chờ socket. Thread đó là một
trong bốn event loop, và trong lúc chờ nó không phục vụ hàng trăm
socket khác mà nó sở hữu. Request của chúng khựng lại --- không thất
bại, \emph{khựng} --- cho tới khi truy vấn trả về. Dưới một cơ sở dữ
liệu chậm, cả bốn vòng lặp bị ghim trong vài giây và gateway nhận kết
nối mà nó không xử lý nổi: mọi route, mọi user, một triệu chứng. Điều
tương tự đúng với \code{.block()} trên một \code{Mono},
\code{Thread.sleep}, một lời gọi \code{RestTemplate} đồng bộ, hay đọc
file. Phương án reactive dời công việc ra khỏi vòng lặp và
\emph{ghép} nó vào chuỗi:

\begin{lstlisting}[language=Java]
return Mono.fromCallable(() -> deny.isRevoked(claims.getId()))
        .subscribeOn(Schedulers.boundedElastic())   // off the loop
        .flatMap(revoked -> revoked
                ? unauthorizedResponse(exchange, "Token revoked")
                : chain.filter(exchange.mutate().request(mutated).build()));
\end{lstlisting}

Tốt hơn nữa là một client không chặn (Redis reactive, R2DBC) để không
thread nào phải chờ cả. Quy tắc cho bất cứ thứ gì trên đường đi của
gateway: \emph{nếu phương thức của thư viện trả về một giá trị thay vì
một} \code{Mono}\emph{/}\code{Flux}\emph{, nó không thuộc về filter.}
Agent BlockHound của Reactor biến quy tắc đó thành một test thất bại.

\subsection{Predicate, thứ tự, và phép tính \code{StripPrefix}}

Route được thử theo thứ tự khai báo (tất cả có \code{order} mặc định
là 0, và phép sắp xếp là ổn định); predicate khớp đầu tiên thắng. Bảng
của dự án không có chồng lấn, nhưng một route bắt-tất-cả
\code{Path=/api/**} trong tương lai đặt \emph{phía trên} các route cụ
thể sẽ nuốt chúng trong im lặng --- đặt route bắt-tất-cả ở cuối hoặc
cho nó \code{order} cao. Các filter gắn vào một route được đánh số tuần
tự và chạy \emph{sau} các global filter có order thấp hơn;
\code{getOrder()} của filter JWT bằng $-1$ đặt nó trước mọi filter của
route, và đó là điều cho phép rate limiter của route (bên dưới) đọc
header \code{X-User-Id} mà filter JWT vừa ghi.

\begin{pitfall}[đoạn path lặp đôi]
Triệu chứng: frontend gọi \code{/api/courses/courses/5} và chạy được,
nhưng gõ tay \code{/api/courses/5} thì course-service trả 404. Nguyên
nhân: \code{StripPrefix=2} bỏ đúng hai đoạn, \code{/api} và
\code{/courses}, nên \code{/api/courses/5} tới nơi thành \code{/5} ---
và không controller nào được ánh xạ ở đó. \code{courses} đầu tiên chọn
\emph{route}; \code{courses} thứ hai sống sót qua phép bỏ và khớp
\code{@RequestMapping("/courses")}. Frontend cố ý lặp đoạn đó. Thiết
kế sạch hơn là \code{StripPrefix=1} với controller giữ tiền tố tài
nguyên của mình; dự án chọn giữ cho service không biết tới \code{/api}
với cái giá là URL lặp đôi.
\end{pitfall}

\section{CORS ở biên}

\code{CorsConfig} đăng ký một \code{CorsWebFilter} cho \code{/**}:
origin \code{localhost:3000} và \code{5173}, sáu phương thức, header
bất kỳ, cho phép credentials. Ba điều cần hiểu về chỗ nó sống.

\emph{Vì sao ở đây mà không ở từng service.} CORS là quy tắc của trình
duyệt về \emph{origin}, và origin mà trình duyệt nói chuyện là gateway.
Một service trả lời \code{Access-Control-Allow-Origin} là trả lời một
câu hỏi không ai hỏi nó; và nếu hai service bất đồng, trình duyệt sẽ
thấy chính sách không nhất quán cho một API. Mối quan tâm của biên đặt
ở biên.

\emph{Preflight tốn gì.} Với bất kỳ request nào không ``đơn giản'' ---
một \code{POST} JSON, bất cứ thứ gì có header \code{Authorization} ---
trình duyệt gửi trước một request \code{OPTIONS} và đợi câu trả lời.
Đó là thêm một vòng đi-về cho mỗi lời gọi, và mọi lời gọi API trong dự
án này đều mang \code{Authorization}. Gateway tự trả lời preflight mà
không đụng tới route nào, rất rẻ, nhưng vòng đi-về đó là người dùng
trả. \code{Access-Control-Max-Age} cho trình duyệt cache câu trả lời;
\code{CorsConfig} không đặt nó, nên Chrome rơi về năm giây. Một dòng,
\code{config.setMaxAge(3600L)}, loại gần hết preflight khỏi một phiên.

\emph{Khi nào nó không bao giờ được hỏi tới.} Trong cụm, frontend và
\code{/api} chung một origin phía sau Traefik (Chương~26), nên không
request nào là xuyên origin và filter đứng không. Dev là nơi duy nhất
nó chạy --- và proxy của Vite có thể khiến ngay cả dev cũng cùng
origin. Vì thế cấu hình CORS ở đây không phải tính năng bảo mật; nó là
tiện nghi phát triển không được mở rộng cẩu thả:
\code{allowCredentials(true)} với origin ký tự đại diện bị trình duyệt
từ chối chính vì nó sẽ cho bất kỳ site nào gọi API với cookie của
người dùng.

\section{Việc của biên, và việc không phải}

\begin{center}
\small
\begin{tabular}{@{}p{0.44\linewidth}p{0.48\linewidth}@{}}
\toprule
Thuộc về biên & Không thuộc \\
\midrule
Xác thực, header danh tính & Quy tắc phân quyền theo từng tài nguyên \\
Rate limit, hạn ngạch & Kiểm tra hợp lệ nghiệp vụ \\
Timeout, retry (lũy đẳng), circuit breaker & Gộp response từ nhiều service \\
CORS, kết thúc TLS, nén & Gọi cơ sở dữ liệu \\
Header truy vết, log truy cập & Biến đổi dữ liệu ngoài header/path \\
\bottomrule
\end{tabular}
\end{center}

Cột bên phải là anti-pattern ``smart pipes'' mà thời ESB đã dạy: khi
gateway biết quy tắc nghiệp vụ, mọi tính năng phải ship ở hai nơi và
đội gateway thành nút thắt của tất cả. Ngoại lệ chính đáng duy nhất
--- ghép nhiều lời gọi service thành một response theo hình dạng màn
hình --- là một service riêng (\emph{backend for frontend}), triển
khai và mở rộng tách khỏi biên.

\section{Ưu và nhược, thẳng thắn}

\begin{itemize}
  \item \textbf{Ưu:} một địa chỉ, một filter JWT, một chính sách CORS;
  route là cấu hình, đổi được mà không cần build lại (Chương~2).
  \item \textbf{Ưu:} mô hình reactive giữ hàng nghìn request đang bay
  trên bốn thread và khởi động trong vài giây.
  \item \textbf{Nhược:} gateway của dự án \emph{không} đặt timeout nào
  cho HTTP client. Mặc định của Reactor Netty là 30 giây để kết nối và
  \emph{không có} timeout phản hồi.
  \item \textbf{Nhược:} không circuit breaker, không rate limiter ---
  client nào cũng gửi nhanh tùy thích, và một service chết được từng
  request phát hiện riêng lẻ.
  \item \textbf{Nhược:} \code{replicas: 1} trong Kubernetes. Thành phần
  mọi request đều đi qua lại là thứ chỉ có một bản.
  \item \textbf{Nhược:} starter nó dùng
  (\code{spring-cloud-starter-gateway}) bị đánh dấu lỗi thời trong
  Spring Cloud 2025.0, thay bằng starter
  \code{gateway-server-webflux}; việc đổi tên hôm nay vô hại và sẽ là
  lỗi build ở một lần nâng cấp tương lai.
\end{itemize}

\section{Cái gì gãy}

\begin{pitfall}[một downstream chậm kéo cả route theo]
Triệu chứng: course-service chậm (một lock, một GC dài), và trong vòng
một phút \emph{mọi} request tới \code{/api/courses/**} thất bại với 503
hoặc treo --- kể cả những request rẻ đáng lẽ thành công. Nguyên nhân:
gateway giữ một connection pool theo từng host, cỡ
\code{max(cores, 8) $\times$ 2} = 16 trên máy này. Không có timeout
phản hồi, mỗi request chậm giữ một kết nối trong pool lâu bằng thời
gian course-service cần; request thứ mười bảy đợi trong hàng
acquire 45 giây rồi thất bại. Mọi user xếp hàng sau người chậm nhất.
Sửa: một timeout phản hồi (Mục~\ref{sec:gw-enhance}) để kết nối kẹt
được thu hồi, và một circuit breaker để route thất bại nhanh khi lỗi
dồn cụm.
\end{pitfall}

\begin{pitfall}[431 Request Header Fields Too Large]
Triệu chứng: một user có nhiều cookie, hoặc một token với claim béo
(Chương~28), nhận ngay 431 từ gateway mà log service không có gì.
Nguyên nhân: Netty giới hạn \emph{khối header} request ở 8\,KB theo mặc
định; gateway từ chối request trước khi bất kỳ filter nào chạy. Sửa:
giữ token mỏng. Tăng \code{server.max-http-request-header-size} là câu
trả lời hạng nhì, vì mỗi chặng phía sau gateway có giới hạn riêng.
\end{pitfall}

\begin{pitfall}[gateway là điểm lỗi đơn]
Triệu chứng: một lần rolling update gateway, một lần reboot node, hoặc
một OOM kill --- và toàn bộ ứng dụng sập, trong khi mọi service phía
sau đều khỏe. Nguyên nhân: một replica. Gateway là phi trạng thái, nên
cách sửa là \code{replicas: 2} cộng một \code{PodDisruptionBudget};
thứ duy nhất làm việc mở rộng khó hơn là trạng thái theo từng instance
như một rate limiter trong bộ nhớ, và đó chính là lý do phần nâng cấp
bên dưới giữ trạng thái đó trong Redis.
\end{pitfall}

\section{Chuyện gì gãy lúc 3 giờ sáng}

Postgres của course-service bắt đầu một đợt autovacuum dài; truy vấn
từng mất 20\,ms giờ mất 4\,s. Đi theo đường đi.

\begin{enumerate}
  \item Frontend bắn sáu request của một dashboard. Bốn cái chạm
  course-service và mất 4\,s mỗi cái. Gateway giữ bốn kết nối trong
  pool suốt 4\,s.
  \item User bấm lại. Hai mươi user, mỗi người bốn request: 80 request
  tranh 16 kết nối. Hàng acquire dài ra; sau 45\,s các request đang đợi
  thất bại với 503.
  \item Request tới auth-service và dictionary-service \emph{không bị
  ảnh hưởng} --- pool là theo host --- cho tới khi access token của user
  hết hạn và \code{/refresh} thành công nhưng dashboard vẫn lỗi, nên họ
  đăng xuất rồi đăng nhập lại, thêm tải.
  \item Không ai bị gọi, vì gateway ``khỏe'': \code{/actuator/health}
  trả 200 và event loop của nó không bị chặn. Sự cố vô hình ở biên và
  chỉ hiện ra dưới dạng độ trễ trong Zipkin (Chương~24).
\end{enumerate}

Với timeout phản hồi 5\,s và circuit breaker trên route, bước~2 trở
thành: các lời gọi chậm bị cắt ở 5\,s và tính là thất bại; sau mười
lần, breaker mở; trong 10\,s tiếp theo mọi request \code{/api/courses}
thất bại dưới một mili giây với một body fallback; pool không bao giờ
cạn; auth và dictionary vẫn chạy; thay đổi trạng thái của breaker là
một metric bạn có thể đặt cảnh báo. Cơ sở dữ liệu vẫn chậm --- khả
năng chịu lỗi không sửa nguyên nhân --- nhưng bán kính thiệt hại là một
route thay vì cả site.

\section{Nâng cấp giải pháp}
\label{sec:gw-enhance}

Ba bổ sung, toàn bộ là cấu hình cộng một bean. Dependency trước
(\code{pom.xml} của gateway):

\begin{lstlisting}[language=yaml]
# spring-cloud-starter-circuitbreaker-reactor-resilience4j
#   -> the reactive CircuitBreaker gateway filter
# spring-boot-starter-data-redis-reactive
#   -> backing store for RequestRateLimiter (non-blocking client)
\end{lstlisting}

Rồi \code{api-gateway.yml} trong config-server:

\begin{lstlisting}[language=yaml]
spring:
  data:
    redis:
      host: ${REDIS_HOST:localhost}            # redis:6379 in compose
  cloud:
    gateway:
      httpclient:
        connect-timeout: 2000                  (*@\co{1}@*)
        response-timeout: 5s
        pool:
          max-connections: 64                  (*@\co{2}@*)
          acquire-timeout: 2000
      default-filters:
        - RemoveRequestHeader=X-User-Id        (*@\co{3}@*)
        - RemoveRequestHeader=X-User-Role
      routes:
        - id: course-service-courses
          uri: lb://course-service
          predicates:
            - Path=/api/courses/**
          filters:
            - StripPrefix=2
            - name: CircuitBreaker             (*@\co{4}@*)
              args:
                name: courseCb
                fallbackUri: forward:/fallback/course
                statusCodes: "500,502,503,504"
            - name: RequestRateLimiter         (*@\co{5}@*)
              args:
                redis-rate-limiter.replenishRate: 20
                redis-rate-limiter.burstCapacity: 40
                redis-rate-limiter.requestedTokens: 1
                key-resolver: "#{@userKeyResolver}"
resilience4j:
  circuitbreaker:
    instances:
      courseCb:
        sliding-window-size: 10
        minimum-number-of-calls: 10
        failure-rate-threshold: 50
        slow-call-duration-threshold: 3s       (*@\co{6}@*)
        slow-call-rate-threshold: 80
        wait-duration-in-open-state: 10s
        permitted-number-of-calls-in-half-open-state: 3
  timelimiter:
    instances:
      courseCb:
        timeout-duration: 5s                   (*@\co{7}@*)
\end{lstlisting}

\begin{description}
  \item[\co{1}] Kết nối trong 2\,s hoặc bỏ cuộc; phản hồi đầy đủ trong
  5\,s hoặc kết nối bị đóng và trả về pool. Có thể ghi đè theo route
  qua \code{metadata} của route cho một endpoint chậm đã biết như tải
  file lên.
  \item[\co{2}] Mười sáu kết nối mỗi host là chật một khi timeout làm
  request ngắn lại; 64 là thoải mái trên máy này mà vẫn có giới hạn.
  \item[\co{3}] Cách sửa của Chương~28: header danh tính chỉ có thể do
  gateway ghi.
  \item[\co{4}] Filter Resilience4j reactive. Thất bại (và các mã trạng
  thái liệt kê) tính vào breaker; khi mở, request được chuyển tiếp tới
  \code{/fallback/course} \emph{bên trong gateway}.
  \item[\co{5}] Token bucket trong Redis: nạp lại 20 token mỗi giây,
  tích tối đa 40, mỗi request tiêu một. Vượt quá thì trả 429 với
  \code{X-RateLimit-Remaining: 0}. Redis chứ không phải bộ nhớ, để hai
  replica gateway dùng chung một bucket cho mỗi user.
  \item[\co{6}] Không có dòng này, lời gọi chậm không phải thất bại ---
  ngưỡng chậm mặc định là 60\,s. Ba giây nghĩa là ``chậm hơn thiết
  kế''.
  \item[\co{7}] Breaker của gateway bọc mỗi lời gọi trong một
  TimeLimiter có mặc định là \textbf{một giây}. Bỏ trống, một lần tải
  lên khỏe mạnh 1,2\,s sẽ làm breaker nhảy. Đặt tường minh và giữ nó
  $\ge$ timeout phản hồi.
\end{description}

Key resolver và endpoint fallback:

\begin{lstlisting}[language=Java]
@Configuration
public class RateLimitConfig {

    /** Bucket per authenticated user; anonymous callers share
     *  a bucket per client IP (public paths: login, refresh). */
    @Bean
    public KeyResolver userKeyResolver() {
        return exchange -> {
            String user = exchange.getRequest().getHeaders()
                    .getFirst(JwtConstants.HEADER_USER_ID);
            if (user != null) {
                return Mono.just("u:" + user);
            }
            var addr = exchange.getRequest().getRemoteAddress();
            return Mono.just("ip:" + (addr == null ? "unknown"
                    : addr.getAddress().getHostAddress()));
        };
    }
}

@RestController
public class FallbackController {

    /** Served by the gateway itself when a breaker is open.
     *  503 + Retry-After tells clients (and browsers) to back off. */
    @GetMapping("/fallback/course")
    public ResponseEntity<Map<String, String>> course() {
        return ResponseEntity.status(HttpStatus.SERVICE_UNAVAILABLE)
                .header(HttpHeaders.RETRY_AFTER, "10")
                .body(Map.of("error", "course-service unavailable",
                             "retryAfterSeconds", "10"));
    }
}
\end{lstlisting}

\begin{decision}[giới hạn tốc độ theo khóa nào]
Theo \emph{user id} là mặc định đúng cho một API đã xác thực: nó khớp
với thứ bạn đang bảo vệ (phần công bằng của một tenant) và không né
được bằng cách đổi IP. Theo \emph{IP} là lựa chọn duy nhất trước khi
xác thực --- login và refresh --- và là lựa chọn đúng ở đó, vì mối đe
dọa là nhồi credential từ một nguồn; nó sai phía sau NAT công ty, nơi
một IP là một nghìn người dùng. Theo \emph{API key} dành cho client máy
và đối tác, nơi key là hợp đồng. Biên production thường xếp chồng hai
lớp: giới hạn thô theo IP trên đường dẫn công khai và giới hạn theo
user trên mọi thứ còn lại.
\end{decision}

\begin{handson}[đập vào bucket]
Với Redis đã thêm vào \code{docker-compose.yml} và gateway khởi động
lại, bắn 50 request với tư cách một user nhanh hơn bucket nạp lại:
\begin{lstlisting}[language=bash]
$ for i in $(seq 1 50); do
    curl -s -o /dev/null -w "%{http_code} " \
      -H "Authorization: Bearer $TOKEN" \
      http://localhost:8080/api/courses/courses
  done; echo
200 200 200 ... 200 429 429 429 429 429 429 429 429 429 429
$ curl -si -H "Authorization: Bearer $TOKEN" \
    http://localhost:8080/api/courses/courses | grep -i ratelimit
X-RateLimit-Remaining: 0
X-RateLimit-Burst-Capacity: 40
X-RateLimit-Replenish-Rate: 20
\end{lstlisting}
Bốn mươi cái thành công (phần burst), rồi 429 cho tới khi một giây
trôi qua và hai mươi token quay lại. Dừng course-service và lặp lại:
sau mười thất bại, cùng lệnh đó trả về JSON fallback dưới một mili
giây, và \code{/actuator/circuitbreakers} cho thấy \code{courseCb} ở
trạng thái \code{OPEN}.
\end{handson}

\section{Ngoài dự án}

Kong, Envoy (và ingress của Istio xây trên nó), nginx, và các gateway
cloud (AWS API Gateway, Apigee) hiện thực cùng bốn ý tưởng --- timeout,
breaker, rate limit, viết lại header --- dưới dạng cấu hình gốc thay vì
filter Spring, với token bucket thường nằm trong bộ nhớ riêng của proxy
trên từng node cộng một kho chung cho giới hạn toàn cục. ``Outlier
detection'' của Envoy là circuit breaker; filter ``local rate limit''
của nó là bucket. Điều bạn học ở đây chuyển thẳng sang: những câu hỏi
người phỏng vấn hỏi về bất kỳ thứ nào trong số đó --- trạng thái nằm ở
đâu, chuyện gì xảy ra khi kho chết, thứ bậc timeout là gì --- có cùng
câu trả lời.

\section{Câu hỏi từ phỏng vấn thật}

\subsection*{Q: Lưu lượng tăng 10$\times$ sáng thứ Hai. Cái gì gãy trước trong gateway, và bạn tinh chỉnh gì?}

Không phải event loop. Bốn vòng lặp làm I/O không chặn xử lý được
nhiều hơn mười lần lưu lượng của dự án này rất xa; CPU trên gateway là
xác minh HMAC và chép header, chỉ vài micro giây. Cái gãy trước là bất
cứ thứ gì \emph{không có giới hạn}: connection pool theo host (mặc định
16 ở đây) đầy ngay khi bất kỳ downstream nào chậm, và không có timeout
phản hồi thì nó không bao giờ rút. Nên thứ tự tinh chỉnh là: timeout
phản hồi và kết nối để kết nối được thu hồi; nâng trần pool lên mức
downstream thực sự phục vụ được (\code{max-connections}, và một
\code{acquire-timeout} ngắn để request xếp hàng thất bại nhanh thay vì
chất đống); circuit breaker theo route để service đang vật lộn bị cắt
bớt thay vì xếp hàng; rồi mở rộng ngang --- hai hoặc ba replica phía
sau Service với HPA theo CPU, việc tầm thường với một gateway phi trạng
thái một khi trạng thái rate limit nằm trong Redis. Giới hạn bộ nhớ
cũng quan trọng: Netty giữ buffer trực tiếp ngoài heap, nên giới hạn
512\,Mi với \code{MaxRAMPercentage=75} chừa 128\,Mi cho chúng; dưới
lượng đệm 10$\times$, đó là nơi một OOM kill sẽ đến từ.

\subsection*{Q: Một downstream chậm và cả site treo. Vì sao, và bạn làm gì?}

Vì site dùng chung thứ đang cạn kiệt. Nếu gateway có một connection
pool theo host, chỉ route của host đó phải suy giảm --- nên ``cả site''
nói lên có thứ khác đang dùng chung: một event loop bị chặn bởi lời
gọi đồng bộ trong filter, một rate limiter toàn cục khóa sai thứ, hoặc
frontend bắn mọi lời gọi trong một \code{Promise.all} nên một phản hồi
chậm giữ cả trang. Sửa theo thứ tự đó: rà filter tìm lời gọi chặn
(BlockHound trong test), cho route chậm một timeout và breaker để nó
thất bại trong mili giây với body fallback, và cho frontend render
trang từng phần. Điểm senior là khả năng chịu lỗi nói về \emph{cô lập}:
mục tiêu không phải course-service không bao giờ chậm mà là sự chậm của
nó không với tới được trang đăng nhập.

\subsection*{Q: Làm sao rate limit theo user xuyên ba replica gateway?}

Trạng thái phải dùng chung, nếu không giới hạn là theo từng replica và
lệch ba lần. \code{RequestRateLimiter} của Spring Cloud Gateway chạy
một token bucket dưới dạng script Lua trong Redis, nguyên tử, khóa theo
bất cứ thứ gì \code{KeyResolver} trả về --- ở đây \code{u:<userId>} từ
header mà filter JWT đã đặt, rơi về IP của client trên đường dẫn công
khai. Các con số là chính sách: \code{replenishRate} là số request mỗi
giây duy trì được, \code{burstCapacity} là thứ một lần tải trang được
tiêu cùng lúc; đặt burst cỡ số lời gọi của một màn hình. Redis thành
một phụ thuộc, nên quyết định chế độ lỗi của nó một cách có chủ ý: mặc
định của filter là \emph{cho qua} khi Redis không với tới được, đúng
với ứng dụng nội bộ và sai với ứng dụng công khai đối mặt lạm dụng. Với
lưu lượng rất cao, một bucket cục bộ theo replica đặt trước bucket dùng
chung hấp thụ phần lớn vòng đi-về tới Redis.

\subsection*{Q: Bạn đặt xác thực ở đâu, và không đặt logic ở đâu?}

Xác thực --- ``token này hợp lệ không và của ai'' --- ở biên, một lần,
vì đó là cùng một câu hỏi cho mọi request và vì đó là nơi credential đi
vào hệ thống. Phân quyền --- ``user này được chấm bài nộp này không''
--- trong service sở hữu bài nộp, vì chỉ nó biết dữ liệu. Gateway có
thể cưỡng chế các quy tắc thô vẫn còn hình dạng request
(``\code{/api/admin/**} cần vai trò ADMIN'') nhưng khoảnh khắc một quy
tắc cần một dòng trong cơ sở dữ liệu, nó thuộc về phía sau. Không gì có
hình dạng nghiệp vụ nằm trong gateway: không gộp, không biến đổi body,
không gọi kho dữ liệu. Phép thử tôi dùng khi review: nếu một thay đổi
quy tắc miền đòi deploy gateway, ranh giới đã bị vượt.

\subsection*{Q: Làm sao deploy gateway mà không rơi request nào?}

Tối thiểu hai replica, một \code{PodDisruptionBudget} với
\code{minAvailable: 1}, và rolling update với \code{maxUnavailable: 0}
để pod mới Ready trước khi pod cũ bị gỡ. Readiness phải thật: probe
nên thất bại khi registry Eureka còn rỗng, nếu không pod mới nhận
request trước khi nó route được đi đâu. Lúc tắt, pod cũ phải ngừng
nhận \emph{và} hoàn tất request đang bay: \code{server.shutdown:
graceful} với \code{spring.lifecycle.timeout-per-shutdown-phase} dài
hơn timeout phản hồi, và một \code{preStop} sleep vài giây để việc gỡ
endpoint của kube-proxy lan truyền trước khi tiến trình nhận SIGTERM.
Khoảng hở cuối đó --- danh sách endpoint và vòng đời pod do hai
controller khác nhau cập nhật --- là thứ làm rơi các request mà người
ta đổ cho ``rolling update''.

\subsection*{Q: Một tích hợp đối tác cần một response ghép từ ba service. Gateway?}

Không. Đó là một backend-for-frontend: một service nhỏ gọi ba service
song song, chịu được một cái lỗi, và định hình response cho client đó.
Nó mở rộng và deploy theo lịch riêng, sở hữu timeout riêng, và giữ
gateway là một proxy. Đặt nó vào gateway thì mọi thay đổi của đối tác
là một lần deploy biên, và event loop của gateway bắt đầu giữ trạng
thái fan-out cho lời gọi chậm nhất trong ba.

\begin{quiz}
\begin{enumerate}
  \item Filter JWT trả về \code{chain.filter(exchange)} mà không đợi
  gì. Giải thích cái gì đã và chưa xảy ra vào lúc dòng đó trả về, và vì
  sao một filter servlet không thể viết theo cách này.
  \item Một đồng nghiệp thêm một lượt tra Redis đồng bộ vào filter JWT
  để kiểm tra danh sách cấm. Dưới một lần Redis khựng 200\,ms, chuyện
  gì xảy ra với request tới các route \emph{khác}, và vì sao tác động
  trên máy này là 25\,\% cho mỗi lời gọi bị khựng?
  \item course-service treo ở 10\,\% request. Với cấu hình gateway hiện
  tại của dự án, bao nhiêu request treo đồng thời sẽ chặn \emph{toàn
  bộ} lưu lượng tới service đó, và thiết lập duy nhất nào thay đổi con
  số ấy?
  \item Filter \code{CircuitBreaker} của gateway được thêm vào một
  route mà không có khối \code{timelimiter}. Một lần tải lên khỏe mạnh
  mất 1,5\,s. Chuyện gì xảy ra, và vì sao?
  \item Vì sao trạng thái của rate limiter nằm trong Redis thay vì bộ
  nhớ của gateway, và đánh đổi là gì khi Redis chết?
  \item Mọi request từ frontend đều mang \code{Authorization}. Điều đó
  hàm ý gì cho preflight CORS khi dev, thiết lập duy nhất nào giảm
  chúng, và vì sao tất cả không áp dụng trong cụm?
\end{enumerate}
\end{quiz}
